\section{Discussion}
\label{sec:discussion}

\paragraph{What generalizes.}
Nothing in Sections~\ref{sec:drift}--\ref{sec:skip} used two dimensions, and
nothing used regular polygons beyond the value of $\kappa$. The argument needs
exactly four things: a metric that is a support function of some convex body
inside the unit ball, a lower bound $\kappa\lVert q\rVert$ on it, an index-affine
motion of the frames, and a renderer that clamps the field at a known guard. Any
family of nested convex level sets under a one-parameter rigid motion admits the
same interval, with $\kappa$ the inradius of the normal hull and
$m=\rad(-\delta)$ the support value along the walk. In $3$D the same construction
covers nested convex shells under a screw motion; $\kappa$ becomes the inradius of
the normal hull of the polyhedron, and Proposition~\ref{prop:closed} becomes a max
over faces rather than edges. We have not implemented this, but the proofs do not
change.

Conversely, the method fails cleanly where the hypotheses fail. A non-convex
generator has no support-function representation, so (P2) and (P4) go and with them
the interval and the closed form; the Lipschitz skip survives, since it needs only
(P1). A motion that is not affine in the index (rings whose rotation
accelerates, say) loses the affine envelopes but keeps a Lipschitz bound if the
motion is Lipschitz in $n$, so Algorithm~\ref{alg:scan} would still run, only
without a certificate.

\paragraph{Why the discrete setting is favorable.}
Why does this problem admit a stronger guarantee than the
continuous analogues in Section~\ref{sec:related}? Range analysis on a spatial
domain must subdivide, because the domain is uncountable and the bound is
conservative; it terminates on a tolerance. Here the domain is $\N_0$, the interval
is exact rather than conservative, and ``enumerate the interval'' is a terminating
algorithm that returns the true minimum, not an approximation of it. Discreteness
converts a tolerance into a proof. We suspect other rendering problems phrased as
``which member of this indexed family is nearest'' (glints~\cite{Yan2014,
Yan2016}, tiled or stamped patterns~\cite{Lagae2010}, layered fibers) have the
same structure available and are being solved with budgets instead.

\paragraph{What the representation provided, besides speed.}
We began with the field formulation as a way to get resolution independence, and
that is what it delivered. What we did not expect is that it would function as a
debugging method. Four shipped bugs were found the same way, by writing the
layer's definition down and watching a unit fail to agree: the missing eikonal
divide (\S\ref{sec:eikonal}; supplemental \S S2), a stroke floor applied in
phase units, a subsampling anchor nobody had noticed was a choice
(\S\ref{sec:results}), and an envelope pivot that counted a hex lattice as one
family when it draws three (\S\ref{sec:envelope}). A representation that names
the quantities forces the units to agree, and most of our bugs were units. It
also says which quantities must never be written twice. Each of our first six
fields carried a gradient derived by hand and transcribed into two languages;
the definition of a layer says that gradient is a consequence of the field, not
an attribute of it, and forward-mode differentiation (\S\ref{sec:exprs}) removes
the transcriptions along with the class of bug they were the only source of.

\paragraph{Where the ink runs out.}
The formulation's real budget is not arithmetic, it is coverage. A layer costs one
distance query, so twelve layers cost twelve; but twelve dense families of ink in
one frame is a black rectangle, and the pattern is gone long before the frame rate
is. This bound binds in exactly one place we wanted to go: displaying two encoded
fields at once (\S\ref{sec:beyond}) needs four families, and four families of
strokes at a pitch fine enough to carry fringes saturates. Hue and pitch separation
help and are not enough. We think the missing piece is a compositing rule for
fields that is not union-of-ink, something that spends contrast rather than
coverage, and we do not have one. Of the open problems in this paper it is the one that matters most, because it
is what stands between a moir\'e and a display with more than one channel.

\paragraph{Limitations.}
Five, in decreasing order of severity. \emph{The scan's reach follows its
gradients.} The selection of \S\ref{sec:selection} is order-free only where
both index gradients are closed-form; a walking family's gradient is a
screen-space difference, whose per-quad noise the reduction reads as slow
high-order characters, so a pair involving one keeps the old order-two floor,
and a lattice's candidate combinations are still an enumeration (the first
three dual rings, supplemental \S S3), so a sufficiently extreme lattice pair
can beat outside them. Both caps are now exceptions that name their cause
rather than a ceiling on the theory.
\emph{The marginal band.} On
$\rad(-\delta)=s$ with $\theta\neq0$ the interval is genuinely infinite and we
neither enumerate it nor prefilter it; we subsample, coherently, and accept up to
$\degLegibleWorst\%$ ink loss in the legible corner of that band
(Section~\ref{sec:limits}). Proposition~\ref{prop:closed} solves the $\theta=0$
slice exactly, which suggests the rotated case is not hopeless: the residual there
is a sinusoid plus a line, and the indices where it can vanish are governed by the
continued-fraction expansion of $\theta/2\pi$: the same arithmetic
\S\ref{sec:selection} made the criterion's, so the tools are already in the
renderer. We consider it the most promising follow-on. \emph{The metric.} We minimize the residual of the inradius metric, not
Euclidean distance to the curve, which thickens strokes near polygon corners by up
to $1/\kappa$; dividing by $\lVert\nabla\rad\rVert$ would fix it, at the cost of
the exact envelopes. \emph{Prefiltering.} We show the wide-interval regime is
saturated but do not integrate the field there; a correct analytic mean over a
walking family is open, and would replace the subsampling fallback with something
principled rather than merely coherent. \emph{Families without a scalar index.}
Four of our thirteen (the radial pencil and the three lattices) have no
phase function in the sense of \S\ref{sec:catalog}: the pencil's index is an angle, so its fringes
are not level sets of a difference of scalars, and a lattice's index is a pair of
integers, so a lattice pair produces a two-dimensional beat that the fringe law
describes only along a chosen direction. They render correctly, and both are
classical enough that the beat is well understood elsewhere
\cite{Amidror2009Theory,Oster1964}; they simply sit outside the theorem, and a
vector-valued index law that covered them would be a genuine extension.

\paragraph{Future work.}
\label{sec:future}
Beyond the marginal band and prefiltering: the interval is computed per pixel from
$r$, $\Lambda$, and $\kappa$, all of which vary smoothly, so a conservative
interval could be computed once per tile and shared, which on a GPU trades
divergence for a slightly wider scan, probably a win at zoom-out given the
divergence measurements in Table~\ref{tab:ablation}. And since the residual's
zero set is a curve in $(p, n)$ space, a screen-space continuation that carried
$n^\star$ from pixel to pixel would reduce the search to a local refinement almost
everywhere, with the interval available as the certificate that reseeds it when the
continuation breaks.

\section{Conclusion}

Moir\'e has been drawn the same way for a century and a half: put down two layers
of marks and let the overlap do the rest. Drawn that way, the fringe is an
accident of sampling: what the overlap happened to produce, rather than what
was asked for. We have argued that a moir\'e layer is better described by two
scalar fields, an index and a distance, and that with those in hand the fringes
stop being an accident. They are the unit level sets of the difference of the
index fields (Theorem~\ref{thm:fringe}), a statement that under its local
hypotheses needs no periodicity, no lattice, and no Fourier transform, and that
we measured holding on every family we tested, including one whose index
field exists only as the answer to a search.

Three things follow. The pattern becomes a document rather than a raster:
resolution-independent, so that one description is a thumbnail, a screen, and a
plate. The fringe becomes predictable before it is drawn, because the local
ratios $\het_k$ of Eq.~\ref{eq:ratiok} say where one will appear and how it
will run, and \emph{which} beat it will be is no longer scanned for but
computed: the visible-beat hierarchy is the continued fraction of the local
pitch ratio, found at any order by per-pixel lattice reduction
(\S\ref{sec:selection}), so the classical order cap is deleted rather than
raised. And the fringe becomes \emph{addressable}:
choosing the index fields so that their difference is a field of interest turns
the moir\'e into a contour plot of that field, which we verified localizes to a
mean of $\contourWorstErr$\,px against ground truth on the hardest of
\contourFieldCount{} fields, written as expressions and differentiated as they
are evaluated so that a field the author invents arrives with the gradient the
renderer requires; a circle-valued field ends fringes at defects whose signed
endings count its winding exactly (\S\ref{sec:defects}).

Making this run at interactive rates cost the paper's other half. When each member
of a family carries its own rigid motion there is no global index field, and the
per-pixel question becomes an inversion over the integers. We showed that one
classical object, the support function of the generating shape, supplies
everything needed to replace the usual fixed-window guess with a certificate: a
Lipschitz modulus, an exact anisotropy floor $\cos(\pi/N)$, an exact drift rate
$\rad(-\delta)$, and, when the per-index rotation vanishes, convexity of the
residual and hence an exact solution for polygons of any side count, including the
marginal case where no finite enumeration can succeed. That is
$\gpuSpeedLo$--$\gpuSpeedHi\times$ faster than the enumeration it replaces
across rotated settings, and in exact agreement with brute force on every
legible setting of the fidelity sweep (\S\ref{sec:limits} measures the strided
band the sweep's cap excludes).

The work is embodied in \Moire{}, a browser tool in which thirteen families,
twelve simultaneous layers, authored fields, and every parameter in this paper
are live under the cursor at $60$\,Hz. Building it is what produced the
mathematics: the eikonal divide, the pixel-relative stroke floor, and the drift
interval were each found by dragging a slider until something looked wrong. A
field formulation is a good way to make a pattern precise, but an instrument
that moves is a good way to find out which parts of it are false.

Two lessons we would carry elsewhere. A bound that is merely loose, not wrong, can
still produce a visibly broken image whenever the fallback for a loose bound is
approximation, so the tightness of a constant is not always a performance
question. And when approximation is unavoidable, its structure matters more than
its magnitude: the same missing ink reads as a thinner pattern or a shattered one
depending only on whether the samples are anchored to the family or to the pixel.
